\section{Forelesning 4: \emph{Strategi og digital transformasjon}}
\label{sec:lec04}
\settermlabel{terms:lec04}

\subsection*{Oversikt}
Forelesning 4 (Ulrik Røhl) gjør et avgjørende begrepsmessig skille mellom
\term{Digital Transformation (DT)} og \term{IT-Enabled Organizational
Transformation (ITEOT)}, basert på \modelref{wessel}{Wessel et al.\ (2021)}. Distinksjonen er
strategisk viktig fordi den bestemmer \emph{hvilken type} digital
investering som faktisk passer virksomheten -- og dermed hvilken anbefaling
som er forsvarlig. Forelesningen avsluttes med Kemp (2024) sin behandling av
AI som strategisk situeringsoppgave.

\subsection{Hovedteori}

\subsubsection*{\modelref{wessel}{Wessel et al.\ (2021)}: De to typene transformasjon}

\begin{definition}{Digital Transformation (DT)}
Digital teknologi er \emph{sentral for å redefinere} virksomhetens value
propositions, noe som fører til en \emph{ny organisasjonsidentitet}.
Virksomheten blir noe annet enn den var.
\end{definition}

\begin{definition}{IT-Enabled Organizational Transformation (ITEOT)}
Digital teknologi \emph{støtter og forsterker en eksisterende} value
proposition og identitet. Virksomheten gjør det samme den alltid har gjort,
men bedre, raskere eller billigere.
\end{definition}

\keypoint{Forskjellen ligger i hva som skjer med \emph{organisasjonsidentiteten}:
endres den (DT) eller forsterkes den (ITEOT)?}

\begin{figure}[H]
\centering
\includegraphics[width=0.9\linewidth]{lec04-digital-transformation-vs-it-enabled.png}
\caption{Wessel et al.\ (2021): DT redefinerer value proposition og fører til
ny organisasjonsidentitet; ITEOT forsterker den eksisterende.}
\end{figure}

\subsubsection*{\modelref{wessel}{Wessels} fulle prosessmodell}
Wessel et al.\ utvider det enkle skillet til en full prosessmodell:
\begin{itemize}
  \item \term{Technological change} (omgivelses- og organisasjonskontekst)
    $\rightarrow$
  \item \term{Transformation agenda} (eksisterende identitet $\rightarrow$ ny
    identity claim) $\rightarrow$
  \item \term{Transformation activity} (digital teknologi $\leftrightarrow$ value
    proposition) $\leftrightarrow$
  \item \term{Imposition \& reconciliation} (work practice change
    $\leftrightarrow$ reconciliation action) $\rightarrow$
  \item Ny eller forsterket \term{organisational identity}.
\end{itemize}
Modellen viser at transformasjon er en \emph{prosess} med agenda, aktiviteter
og reconciliation -- ikke et engangsvedtak.

\begin{figure}[H]
\centering
\includegraphics[width=0.95\linewidth]{lec04-wessel-transformation-full-model.png}
\caption{Wessel et al.\ -- fullstendig prosessmodell for digital
transformasjon.}
\end{figure}

\subsubsection*{Empirisk grunnlag (Wessel et al.\ 2021)}
\begin{itemize}
  \item \textbf{Case Alpha} -- Et sykehus i Frankrike som innfører elektronisk
    pasientjournal. Identiteten forblir \guillemotleft{}sykehus\guillemotright{}. Klassisk
    \emph{IT-enabled transformation}.
  \item \textbf{Case Beta} -- En finsk produksjonsbedrift som beveger seg fra
    maskinvare til programvare-kontroll. Sitatet er sterkt: \guillemotleft{}Vi var
    maskinfabrikanter \ldots\ men nå trenger ikke kontrollprogramvaren
    maskinvare lenger.\guillemotright{} Identiteten endres fundamentalt -- ekte
    \emph{digital transformation}.
\end{itemize}

\subsubsection*{Hvorfor distinksjonen er strategisk avgjørende}
Feiltolkning fører til feilinvestering:
\begin{itemize}
  \item \emph{Tror} man at virksomheten gjennomgår digital transformasjon, kan
    man argumentere for å bygge en full plattformidentitet -- med
    tilsvarende risiko.
  \item \emph{Ser} man korrekt at det er IT-enabled transformasjon, fokuseres
    investeringen på \emph{hvordan} digital teknologi støtter det
    eksisterende verdiforslaget.
\end{itemize}

\subsubsection*{Kemp (2024): Situering av AI}
Kemp argumenterer at strategisk verdi av AI ikke kommer av modellen alene,
men av tre \term{situating activities}:
\begin{itemize}
  \item \term{Grounding} -- å forankre AI i virksomhetens faktiske
    arbeidspraksis og data.
  \item \term{Bounding} -- å avgrense hvor AI brukes (og hvor mennesket
    fortsatt har siste ord).
  \item \term{Recasting} -- å reformulere oppgaver, roller og verdikjeder slik
    at AI faktisk gir strategisk forskjell.
\end{itemize}

\begin{figure}[H]
\centering
\includegraphics[width=0.85\linewidth]{lec04-kemp-situating-ai-grounding-bounding-recasting.png}
\caption{Kemps tre situering-aktiviteter for AI -- Grounding, Bounding,
Recasting.}
\end{figure}

\begin{figure}[H]
\centering
\includegraphics[width=0.95\linewidth]{lec04-kemp-situated-ai-table.png}
\caption{Kemps tabell: de tre situering-aktivitetene sammenlignet på
definisjon, strategisk begrensning som adresseres, dimensjon av menneskelig
agency og eksempler.}
\end{figure}

\subsubsection*{Kobling til posisjonering (\autoref{sec:lec01})}
DT redefinerer hva virksomheten \emph{er} -- en ny strategisk posisjon. ITEOT
forsterker den eksisterende posisjonen. Porters spørsmål (\guillemotleft{}er aktiviteten en
del av et unikt aktivitetssystem?\guillemotright{}) gjelder begge, men svaret ser ulikt ut.

\subsection{Sentrale fagbegreper}\label{terms:lec04}
\begin{itemize}
  \item \term{Digital Transformation} -- redefinering av value proposition og
    identitet.
  \item \term{IT-Enabled Organizational Transformation (ITEOT)} -- forsterkning
    av eksisterende value proposition.
  \item \term{Organizational identity} -- selvforståelsen av \guillemotleft{}hva vi er som
    virksomhet\guillemotright{}.
  \item \term{Value proposition} -- løftet til kunden.
  \item \term{Identity claim} -- den nye identiteten en transformasjonsagenda
    sikter mot.
  \item \term{Situating activities} (Kemp) -- grounding, bounding, recasting.
\end{itemize}

\subsection{Diagnostisk rammeverk}
For å plassere en konkret virksomhet, still:
\begin{enumerate}
  \item Endres value proposition fundamentalt eller forsterkes den eksisterende?
  \item Endres organisasjonsidentiteten -- blir virksomheten \guillemotleft{}noe annet\guillemotright{}?
  \item Er den digitale teknologien \emph{sentral} (uten den finnes ikke
    tilbudet) eller \emph{støttende} (tilbudet finnes uten, men forbedres)?
\end{enumerate}

\subsection{Eksempler og caser}
\begin{eksempelboks}{Coop: ITEOT, ikke DT}
Coop gjennomgår \emph{IT-enabled transformation}. Appen støtter og
forsterker den butikkbaserte forretningsmodellen og medlemsidentiteten. Coop
redefinerer seg ikke som digital aktør -- bevis: lukkingen av
Coop.dk-nettbutikken i januar 2025.
\end{eksempelboks}

\begin{eksempelboks}{Netflix: DT}
Klassisk DT: gikk fra DVD-utleie til streaming-/produksjonsselskap.
Identiteten endret seg fra ``vi sender ut DVD-er'' til ``vi er et
produksjons- og distribusjonsselskap''.
\end{eksempelboks}

\begin{eksempelboks}{Finsk maskinfabrikant (Beta)}
DT i Wessel et al.: fra maskinvare til programvarekontroll. Den fysiske
maskinvaren ble valgfri -- virksomhetens identitet endret seg fundamentalt.
\end{eksempelboks}

\begin{eksempelboks}{Banker som lanserer mobilbank}
ITEOT så lenge identiteten forblir \guillemotleft{}bank\guillemotright{}; ville blitt DT hvis banken
redefinerte seg som plattformaktør i et bredere finans- og dataøkosystem.
\end{eksempelboks}

\subsection{Eksamensrelevans}
\begin{itemize}
  \item Kunne definere begge typer presist (Wessel et al.\ 2021).
  \item Bruke distinksjonen til å diagnostisere konkrete virksomheter.
  \item Forsvare \emph{hvorfor} distinksjonen er strategisk viktig (forhindrer
    feilinvestering).
  \item Kombinere med Weill \& Woerners DBM (\autoref{sec:lec03}): Coop som
    \emph{omnichannel} + ITEOT $\rightarrow$ digital investering bør styrke
    butikkøkonomien, ikke bygge en separat digital identitet.
  \item Bruke Kemps situating-rammeverk hvis casen involverer AI -- og
    påpeke at strategisk AI-verdi krever organisatorisk arbeid, ikke bare
    teknologi.
\end{itemize}
